\section{EDC / eCRF}

% SLIDE ID: sec02-goal
\BlockGoalFrame{\SlideID{sec02-goal}}{
  \item zrozumieć, czym różni się EDC od samego formularza eCRF,
  \item zobaczyć, jak projekt formularza wpływa na jakość danych,
  \item przygotować się do części demonstracyjnej i warsztatowej.
}

% SLIDE ID: sec02-core
\begin{frame}{EDC i eCRF: punkt centralny pracy z danymi}
  \SlideID{sec02-core}
  \begin{itemize}
    \item EDC służy do elektronicznego zbierania danych badania.
    \item eCRF jest interfejsem i strukturą formularzy, przez które dane trafiają do systemu.
    \item Jakość projektu eCRF wpływa bezpośrednio na jakość danych, monitoring i analizę.
  \end{itemize}
\end{frame}

% SLIDE ID: sec02-scope
\begin{frame}{Zakres sekcji}
  \SlideID{sec02-scope}
  \begin{itemize}
    \item logika formularzy i wizyt,
    \item walidacje i query management,
    \item role użytkowników i audit trail.
  \end{itemize}
\end{frame}

% SLIDE ID: sec02-definition
\begin{frame}{EDC w praktyce}
  \SlideID{sec02-definition}
  \begin{DefinitionBox}
    EDC to system zarządzający strukturą danych badania, przepływem wpisów, walidacjami i kontrolą zmian, a nie tylko „elektroniczna karta obserwacji”.
  \end{DefinitionBox}
\end{frame}
